\subsection{Separating proper actor images from faithful eleven-element sources}

\begin{theorem}[Proved: exact image split for the original $(11,3,5)$ family]
\label{CP2:A-eleven-three-five-faithful-image-split}
Consider a two-object collage with source monoid $M$ of order eleven,
opposite corner $R_3$, and a five-element connecting biset $S$.  Suppose
the minimal ideal of $M$ is a rectangular $3\times2$ band and the collage
has the homotopy type of $\mathbb CP^2$.  Let $q:M\twoheadrightarrow H$
be the image of the two-sided action of $M$ on this band.  Exactly one
of the following two branches applies.
\begin{enumerate}
\item The actor image is proper.  Its order is between seven and ten,
and it is one of the $1+5+16+29=51$ literal faithful hulls in
Proposition~\ref{CP2:A-three-by-two-faithful-hulls-through11}.  Put
$U=S\otimes_M H$.  Then $1\leq |U|\leq5$, and the actual quotient
collage has $H_1(-;\mathbb Z)=0$ and $H_2(-;\mathbb Z)=\mathbb Z$.
Its mod-two degree-two generator must have nonzero square.
Consequently it suffices to check the complete no-fixed-state acts
on two, three, four, and five states for these fifty-one hulls,
retaining every candidate whose mod-two square is nonzero.
\item The actor image has order eleven.  The map $q$ is a monoid
isomorphism, and the induced biset $S\otimes_M H$ is naturally
isomorphic to $S$.  This is precisely the original five-state
collage, up to isomorphism.  Its source is one of the remaining fifty
literal faithful hulls.  This branch may use the ordinary
$\mathbb CP^2$ obstructions, including positive third homology and
an actual proper equivariant biset retraction together with the
proved exclusion of all categories with at most eighteen arrows.
\end{enumerate}
Neither branch requires a census of all monoids of order eleven.
\end{theorem}

\begin{proof}
For a rectangular band, its six elements induce six distinct constant
row/column actor pairs.  Thus the actor quotient is bijective on the
minimal band.  Its identity is distinct from these six actors, and
$7\leq|H|\leq11$.  The complete actor-hull proposition applies because
the relative cup-product lemma and the contractibility of $BR_3$
force the original degree-two generator to have nonzero source
restriction.  In particular the source has
nonzero rational second homology, which supplies the nonzero balanced
functional used in that theorem.

The integral comparison of
Theorem~\ref{CP2:A-band-quotient-integral-low-degree} gives an actual
map of collages inducing integral $H_1$ and $H_2$ isomorphisms.
The set $U$ is a quotient of the nonempty set $S$, so it has between
one and five elements.  Universal coefficients give an isomorphism
on $H^2(-;\mathbb F_2)$.  If the target generator had zero square,
naturality would give zero square for the generator on
$\mathbb CP^2$, a contradiction.  A fixed state on the source side
has unconditional zero products by
Lemma~\ref{CP2:A-hull-square-fixed-state}; this also handles every
one-state biset.  The remaining small-state band forms are exactly
those in Lemma~\ref{CP2:A-three-by-two-small-state-band-forms}.
No assertion about the target's higher homology is needed or made
in this proper-image branch.  In particular, positive target
$H_3$ and a target proper-retract obstruction to being
$\mathbb CP^2$ cannot be used to exclude its square frontier.

If $|H|=11$, surjectivity and the common finite order make $q$
bijective.  A bijective monoid homomorphism is an isomorphism.
Explicitly, the induced-biset isomorphism is
\[
 S\otimes_M H\longrightarrow S,
 \qquad [s,h]\longmapsto s\,q^{-1}(h),
\]
with inverse $s\mapsto[s,1]$.  It respects the left $R_3$ action
and the transported right $H$ action.  Hence the entire quotient
collage is isomorphic to the original one and has exactly five
connecting states.  The usual homological obstructions and the
actual retraction lemma therefore apply to this branch directly.
A proper biset image has at most four states and at most eighteen
arrows; the global eighteen-arrow exclusion rules it out if the
retraction lemma would force it to remain $\mathbb CP^2$.

The ordinary faithful-hull counts through orders seven to eleven
are $1,5,16,29,50$, with the full literal sets and conjugacy scope
independently verified.  This proves the stated $51+50$ partition.
\end{proof}

\begin{remark}[Proved scope; numerical frontiers remain separate]
The square-only controls on eleven-element hulls with three states
found twenty-one nonzero-square records.  These are preserved as
explicit frontiers of the stronger, unneeded all-hull question.
They are not induced three-state quotients in the equal-order
branch of the theorem: an isomorphism preserves the original
five-element state set.  Their presence changes no completed bound.
The two finite branches above require their own closed scope and
literal-certificate joins before an exclusion can be asserted.
\end{remark}
